\documentclass[12pt]{article}
\usepackage{cmap}
\usepackage[T2A]{fontenc}
\usepackage[utf8]{inputenc}
\usepackage[english,russian]{babel}
\usepackage{geometry}
\usepackage{titlesec}
\usepackage{tocloft}
\usepackage{longtable}
\usepackage{array}
\usepackage{amsmath}
\setlength{\emergencystretch}{1em}
\geometry{a4paper, left=2.5cm, right=2.5cm, top=2.5cm, bottom=2.5cm}
\pagenumbering{gobble}

\begin{document}

% =================== ТИТУЛЬНЫЙ ЛИСТ ===================
\begin{titlepage}
    \centering
    \vspace*{\fill}
    
    {\Huge\bfseries Спецификация OLAP базы данных для домоуправления}
    
    \vspace{0.5cm}
    
    {\Large Тимофей Семисоха}
    
    \vspace*{\fill}
\end{titlepage}

\tableofcontents

\newpage
\pagenumbering{arabic}
\setcounter{page}{2}

\section{Назначение аналитической системы}
Система предназначена для анализа начислений, потребления и оплаты коммунальных услуг в жилых квартирах. Она позволяет сопоставлять финансовые показатели по времени, территории, видам услуг и характеристикам квартир, а также оценивать влияние отсутствующих показаний приборов учёта на полноту расчётов.

Аналитическая модель построена по схеме «звезда»: центральная таблица фактов \texttt{olap.fact\_billing} связана с измерениями даты, квартиры, услуги и качества данных. В отличие от транзакционной системы, в данной модели изучаются уже сформированные начисления и платежи, а не операции регистрации жильцов или приёма платежей.

\subsection{Временной охват и детализация}
Расчётный период: с января 2024 года по июнь 2025 года включительно (18 месяцев). Дата контрольного снимка для анализа задолженности - 31 июля 2025 года.

Одна строка таблицы фактов соответствует \textbf{одной квартире, одной коммунальной услуге и одному расчётному месяцу}. Дополнительно в той же таблице хранятся служебные агрегированные строки задолженности. В модели представлены шесть видов услуг: холодное и горячее водоснабжение, электроснабжение, отопление, уборка домовой территории и обслуживание домофона.

\section{Аналитические вопросы}

\subsection{Q1. Начисления по городам и услугам}
\textbf{Вопрос.} Сколько начислено за коммунальные услуги и сколько расчётных записей не удалось рассчитать из-за отсутствия показаний?

\textbf{Показатели:} количество строк начислений, количество строк с неизвестными показаниями, сумма известных начислений; дополнительно - сумма начислений только для групп, где нет ни одной нерассчитанной строки.

\textbf{Группировка:} расчётный месяц, город, вид услуги. \textbf{Период:} январь 2024 - июнь 2025 года.

\subsection{Q2. Годовое изменение начислений}
\textbf{Вопрос.} Как изменились начисления по отношению к тому же месяцу прошлого года?

\textbf{Показатели:} количество сопоставленных записей, суммы начислений за текущий и соответствующий прошлогодний месяцы, процент изменения. Сравниваются только пары записей одной квартиры и услуги с известными начислениями в обоих периодах.

\textbf{Группировка:} месяц текущего года, город, вид услуги. \textbf{Период:} январь--июнь 2025 года относительно января--июня 2024 года.

\subsection{Q3. Удельное потребление электроэнергии}
\textbf{Вопрос.} Сколько электроэнергии потребляют квартиры в расчёте на квадратный метр и как показатель зависит от класса энергоэффективности?

\textbf{Показатели:} количество записей с известным потреблением, суммарное потребление в кВт·ч, суммарная площадь соответствующих квартир и отношение потребления к площади (кВт·ч/м$^2$). Площадь учитывается по записям рассматриваемого месяца.

\textbf{Группировка:} месяц, город, район, класс энергоэффективности из JSONB-поля квартиры. \textbf{Период:} январь 2024 - июнь 2025 года. Учитываются только услуги с единицей измерения кВт·ч и известными показаниями.

\subsection{Q4. Задолженность по сроку просрочки}
\textbf{Вопрос.} Каков размер известной непогашенной задолженности и насколько давно истёк срок оплаты?

\textbf{Показатели:} количество записей с положительным долгом и сумма непогашенной задолженности. Долг определяется как превышение начисленной суммы над оплаченной; записи с неизвестным начислением исключаются.

\textbf{Группировка:} дата снимка, город, вид услуги, интервал просрочки: до 30 дней, 31--90 дней, более 90 дней. \textbf{Период:} состояние на 31 июля 2025 года по начислениям за январь 2024 - июнь 2025 года. Запрос выполняется по предварительно рассчитанным агрегатам через \texttt{olap.v\_aging\_debt}, без повторного расчёта задолженности.

\subsection{Q5. Собираемость платежей}
\textbf{Вопрос.} Какая доля начисленных сумм фактически покрыта платежами и какая доля оплачена своевременно?

\textbf{Показатели:} количество строк с известными начислениями, сумма начислений, покрытая платежами часть начислений, проценты собираемости и своевременной оплаты, сумма переплат. При расчёте покрытия и своевременной оплаты платежи ограничиваются начисленной суммой, поэтому переплата учитывается отдельно.

\textbf{Группировка:} расчётный месяц, город, вид услуги. \textbf{Период:} январь 2024 - июнь 2025 года.

\subsection{Q6. Начисления и долг по характеристикам квартир}
\textbf{Вопрос.} Каковы начисления и задолженность для квартир с классом энергоэффективности C, в которых проводился ремонт?

\textbf{Показатели:} количество расчётных строк, количество записей с неизвестными показаниями, сумма известных начислений и сумма известного долга. Подмножество квартир отбирается по признакам \texttt{energy\_class = C} и \texttt{renovated = true} внутри JSONB-поля.

\textbf{Группировка:} год, полугодие, город, вид услуги. \textbf{Период:} первое и второе полугодия 2024 года, первое полугодие 2025 года.

\subsection{Q7. Потребление воды в квартирах с риском протечек}
\textbf{Вопрос.} Каковы объём потребления и начисления за холодное водоснабжение в квартирах, где есть риск протечек?

\textbf{Показатели:} количество строк начислений, количество записей с неизвестными показаниями, суммарное известное потребление воды в м$^3$ и сумма известных начислений. Квартиры отбираются одновременно полнотекстовым поиском по техническому описанию и проверкой наличия \texttt{pipe\_risk} в массиве тегов.

\textbf{Группировка:} расчётный месяц, город, вид услуги (холодное водоснабжение). \textbf{Период:} январь 2024 - июнь 2025 года.

\subsection{Q8. Пропуски показаний приборов учёта}
\textbf{Вопрос.} Какова доля расчётных строк, для которых отсутствуют показания счётчиков, и сколько денег оплачено по таким строкам при неизвестной сумме начисления?

\textbf{Показатели:} число строк по услугам со счётчиковым способом расчёта, число строк с пропущенными показаниями, их доля в процентах, сумма платежей по строкам с неизвестным начислением.

\textbf{Группировка:} расчётный месяц, город, вид услуги. \textbf{Период:} январь 2024 - июнь 2025 года. В знаменатель доли включаются все строки счётчиковых услуг, а не только строки с пропусками.

\section{Выбор фактов и измерений}
\subsection{Таблица фактов}
\texttt{olap.fact\_billing} хранит количественные показатели по квартире, услуге и месяцу. Её состав определяется необходимостью анализировать одновременно потребление, начисления, платежи и качество исходных данных.

\begin{itemize}
    \item \texttt{quantity} - объём услуги в её единицах измерения; используется для анализа потребления.
    \item \texttt{charged\_amount} - сумма начисления; применяется в отчётах о начислениях, динамике и собираемости.
    \item \texttt{paid\_amount} - фактически оплаченная сумма, включая возможную переплату.
    \item \texttt{paid\_on\_time\_amount} - сумма своевременной оплаты; необходима для анализа платёжной дисциплины.
    \item \texttt{billing\_month} - расчётный месяц; по нему определяется срок оплаты (10-е число следующего месяца).
    \item \texttt{apartment\_key}, \texttt{service\_key}, \texttt{quality\_key} - ключи связей с измерениями квартиры, услуги и качества данных.
\end{itemize}

Если показания счётчика отсутствуют, объём и начисленная сумма могут быть \texttt{NULL}. В таблице срок оплаты не хранится: он однозначно определяется по \texttt{billing\_month}. Дата контрольного снимка также не повторяется в фактах, а хранится в \texttt{olap.dataset}.

\subsection{Предварительно рассчитанные строки задолженности}
Дополнительно в \texttt{fact\_billing} сохраняются агрегированные строки с детализацией \textbf{год $\times$ город $\times$ услуга $\times$ интервал просрочки}. Для каждого сочетания заранее подсчитываются общая сумма известного положительного долга и число расчётных записей с долгом. Интервалы просрочки: до 30 дней, 31--90 дней и более 90 дней на 31 июля 2025 года.

В агрегированных строках \texttt{charged\_amount} содержит сумму долга, а \texttt{quantity} --- количество записей с положительным долгом; \texttt{paid\_amount} и \texttt{paid\_on\_time\_amount} равны нулю. Для отличия от квартир используются отрицательные \texttt{apartment\_key}, ссылающиеся на служебные элементы \texttt{dim\_apartment}; интервал просрочки хранится в их JSONB-свойстве \texttt{aging\_bucket}. Представление \texttt{olap.v\_billing} возвращает только детальные строки, а \texttt{olap.v\_aging\_debt} - агрегированные. Это исключает двойной подсчёт и ускоряет Q4 и годовые сводки по задолженности.


\subsection{Таблицы измерений}
\begin{itemize}
    \item \textbf{Дата - \texttt{olap.dim\_date}.} Содержит 18 первых чисел расчётных месяцев и поля \texttt{date\_key}, \texttt{year\_no}, \texttt{month\_no}. Обеспечивает группировки по месяцам и полугодиям и сопоставление периодов.
    \item \textbf{Квартира - \texttt{olap.dim\_apartment}.} Содержит ключ квартиры, \texttt{building\_id}, город, район, площадь, техническое описание, JSONB-свойства и массив тегов. Служебные элементы с отрицательными ключами кодируют город и интервал просрочки для агрегатов.
    \item \textbf{Услуга - \texttt{olap.dim\_service}.} Определяет шесть услуг, единицу измерения и способ расчёта. Числовые коды расшифровываются представлением \texttt{olap.v\_service}.
    \item \textbf{Качество - \texttt{olap.dim\_quality}.} Различает известные и отсутствующие показания. Необходимо, чтобы отделять неполноту исходных данных от нулевых начислений. Текстовая расшифровка доступна через \texttt{olap.v\_quality}.
\end{itemize}

\section{Целевой объём и наполнение данных}
\subsection{Планируемый объём}
Для наполнения базы задаются следующие параметры:
\begin{itemize}
    \item квартир - 1\,000\,000;
    \item видов коммунальных услуг - 6;
    \item расчётных месяцев - 18 (январь 2024 - июнь 2025 года);
    \item объём детальной части таблицы фактов - \textbf{108\,000\,000 строк};
    \item агрегированных строк задолженности - \textbf{240} (10 городов $\times$ 6 услуг $\times$ 4 сочетания года и интервала просрочки).
\end{itemize}

\[
1\,000\,000\ \text{квартир}\times 6\ \text{услуг}\times 18\ \text{месяцев}
=108\,000\,000.
\]

\subsection{Способ наполнения}
Выбран \textbf{PostgreSQL генератор}: \texttt{olap.load\_data()} с циклами по месяцам и диапазонам квартир. Данные формируются через \texttt{INSERT ... SELECT}, функцией \texttt{generate\_series} с псевдослучайными значениями.

Сначала создаются и заполняются измерения. Затем для каждого из 18 месяцев формируются записи по всем сочетаниям квартир и шести услуг.

\section{Трассировка аналитических вопросов}
Ниже указаны таблица фактов и измерения, участвующие в запросах. \texttt{v\_billing} - представление детальных записей на основе \texttt{fact\_billing}; \texttt{v\_aging\_debt} - представление предварительно рассчитанных агрегатов задолженности.

\begingroup
\small
\renewcommand{\arraystretch}{1.22}
\begin{longtable}{|>{\raggedright\arraybackslash}p{0.8cm}|>{\raggedright\arraybackslash}p{4.2cm}|>{\raggedright\arraybackslash}p{8.2cm}|}
\hline
\textbf{№} & \textbf{Таблица фактов / источник} & \textbf{Используемые измерения} \\
\hline
\endfirsthead
\hline
\textbf{№} & \textbf{Таблица фактов / источник} & \textbf{Используемые измерения} \\
\hline
\endhead
Q1 & \texttt{v\_billing} & \texttt{dim\_date}, \texttt{dim\_apartment}, \texttt{dim\_service} \\
\hline
Q2 & \texttt{v\_billing} (дважды: текущий и прошлогодний месяцы) & \texttt{dim\_date}, \texttt{dim\_apartment}, \texttt{dim\_service} \\
\hline
Q3 & \texttt{v\_billing} & \texttt{dim\_date}, \texttt{dim\_apartment}, \texttt{dim\_service} \\
\hline
Q4 & \texttt{v\_aging\_debt} (агрегаты) & \texttt{dim\_apartment}, \texttt{dim\_service}, \texttt{dataset} (дата снимка) \\
\hline
Q5 & \texttt{v\_billing} & \texttt{dim\_date}, \texttt{dim\_apartment}, \texttt{dim\_service} \\
\hline
Q6 & \texttt{v\_billing} & \texttt{dim\_date}, \texttt{dim\_apartment} (JSONB), \texttt{dim\_service} \\
\hline
Q7 & \texttt{v\_billing} & \texttt{dim\_date}, \texttt{dim\_apartment}, \texttt{dim\_service} \\
\hline
Q8 & \texttt{v\_billing} & \texttt{dim\_date}, \texttt{dim\_apartment}, \texttt{dim\_service}, \texttt{dim\_quality} \\
\hline
\end{longtable}
\endgroup

В Q1--Q3 и Q5--Q7 статус качества определяется по ключу \texttt{quality\_key} таблицы фактов. В Q8 дополнительно используется соединение с \texttt{dim\_quality} для проверки признака \texttt{missing\_reading}. 
В Q4 используются только агрегированные строки; их значения заранее вычислены по известным начислениям с положительным долгом.

\end{document}
